\subsection{El problema}

\texttt{FoodSearchProblem} ya viene completamente implementado en \texttt{searchAgents.py}
(\texttt{getStartState}, \texttt{isGoalState}, \texttt{getSuccessors} y \texttt{getCostOfActions}):
no hubo que escribir código nuevo en esta actividad. El objetivo es entender el problema y
establecer una línea base antes de diseñar heurísticas en la Actividad 11.

El estado ya no es solo la posición de Pac-Man: es el par $s = (p, F)$, donde $p$ es la posición
$(x,y)$ y $F$ es una \texttt{Grid} (matriz booleana) con el estado --presente o consumido-- de cada
alimento del laberinto:

\begin{verbatim}
def isGoalState(self, state):
    return state[1].count() == 0
\end{verbatim}

La meta se alcanza cuando ya no queda ningún alimento por consumir (\texttt{foodGrid.count() == 0}),
sin importar en qué posición esté Pac-Man.

\subsection{Por qué crece tanto el espacio de estados}

Cada alimento puede estar en dos condiciones: presente o consumido. Si el laberinto tiene $F$
alimentos, el número de configuraciones posibles del \texttt{Grid} de comida es, conceptualmente,
$2^F$ --y el estado completo es el producto de eso por el número de posiciones posibles de
Pac-Man--. Esto es exactamente lo que dice la guía, y lo confirmamos experimentalmente: intentamos
correr \texttt{uniformCostSearch} sobre \texttt{smallClassic} (55 alimentos) con un límite de 45
segundos y \textbf{no terminó}. $2^{55}$ es un número astronómicamente más grande que cualquier
espacio de estados manejado en las Actividades 2-9 (donde el estado era solo una posición, o una
posición más 4 bits de esquinas visitadas -- $2^4 = 16$ combinaciones como máximo). Por esta razón,
esta actividad y la 11 se trabajan sobre \texttt{tinySearch} (1 alimento) y \texttt{testClassic}
(8 alimentos), donde UCS puro sigue siendo viable como línea base.

\subsection{Procedimiento y resultados}

Se ejecutaron UCS y A* con \texttt{nullHeuristic} ($h(n)=0$) sobre \texttt{tinySearch} y
\texttt{testClassic} (ver \texttt{experimentos/actividad10\_food\_baseline.py}), replicando la
misma verificación de la Actividad 4 (con $h(n)=0$, A* debe coincidir exactamente con UCS) pero
ahora sobre \texttt{FoodSearchProblem}.

\begin{table}[H]
\centering
\caption{Línea base de UCS y A*+h(n)=0 sobre FoodSearchProblem}
\label{tab:act10-baseline}
\begin{tabular}{lrrrr}
\toprule
\textbf{Layout (alimentos)} & \textbf{Método} & \textbf{Costo} & \textbf{Expandidos} & \textbf{Tiempo (s)} \\
\midrule
tinySearch (1)   & UCS       & 8  & 16   & 0.000393 \\
tinySearch (1)   & A*+h(n)=0 & 8  & 16   & 0.000303 \\
testClassic (8)  & UCS       & 16 & 2598 & 0.063502 \\
testClassic (8)  & A*+h(n)=0 & 16 & 2598 & 0.061340 \\
\bottomrule
\end{tabular}
\end{table}

En ambos layouts, UCS y A*+h(n)=0 encuentran el mismo costo óptimo y expanden exactamente el mismo
número de nodos --la misma consecuencia algebraica de $f(n) = g(n) + 0 = g(n)$ ya explicada en la
Actividad 4--. Nótese el salto de 16 a 2598 nodos expandidos al pasar de 1 a solo 8 alimentos: es
la primera evidencia numérica, dentro de un caso todavía manejable, de la explosión combinatoria que
describe la guía, y es exactamente la motivación para diseñar una heurística no trivial en la
Actividad 11 --sin ella, layouts con más alimentos que \texttt{testClassic} dejan de ser viables,
como mostró el intento fallido sobre \texttt{smallClassic}.
